\documentclass{article}
\usepackage[dblblindworkshop]{neurips_2026}
\workshoptitle{Muslims in ML}
\usepackage[utf8]{inputenc}
\usepackage[T1]{fontenc}
\usepackage{hyperref}
\usepackage{url}
\usepackage{booktabs}
\usepackage{amsmath,amsfonts}
\usepackage{microtype}

\title{Quran Recitation Event Annotation:\\A Competition with Private Human Evaluation}
\author{Anonymous Author(s)}
\begin{document}
\maketitle

\begin{abstract}
A Quran transcript checker must distinguish unresolved mistakes from harmless repetition, successful repairs and accepted spelling differences. We propose a competition that supplies ayah-split transcripts and references, leaves method development to entrants, and scores only their final event predictions against private human annotations. The prepared training release contains 127 unlabelled recording cases and 888 input units, including 44 newly collected cases; a 23-recording review sample comes entirely from training. The private evaluation set contains 100 human-approved cases, 348 units and 162 events across ten labels. An executable label-and-span evaluator and evaluated baselines accompany the task. Agents are permitted, not required, and intermediate iterations are not ranked.
\end{abstract}

\section{Task and community value}
Quran-memorization applications compare noisy ASR transcripts with a fixed public reference. A difference can be an unresolved mistake, a repeated correct phrase, a wrong attempt followed by a repair, or an accepted written variant. Treating every difference as a mistake produces misleading feedback. We evaluate transcript interpretation, not acoustic correctness, pronunciation or tajwid: a possible ASR origin does not excuse a textual mismatch.

For each supplied unit, a method returns events with a combined label and two word spans, one in the original transcript and one in the reference. Spans are zero-based and half-open; missing material uses an empty transcript anchor, extra material an empty reference anchor. No events means clean. Ayah IDs and splits are provided, so identification errors are not scored. Opening formulas are separate units with an empty reference and must also be classified.

Labels are \texttt{substitution\_mistake}, \texttt{omission\_mistake}, \texttt{insertion\_mistake}, \texttt{repetition\_benign}, \texttt{substitution\_corrected}, \texttt{omission\_corrected}, \texttt{letters\_benign}, \texttt{spelling\_benign}, \texttt{isti3adha\_benign} and \texttt{basmala\_benign}. Multiple events may occur in one unit. Adjacent replacements are grouped; repetitions and repairs span both attempts. Ordinary comparison normalization creates no event; contextual spelling equivalences require an explicit rubric rule.

\section{Prepared data and separation}
\textbf{Training.} The frozen text-only release has 127 recording cases: 834 ayah units and 54 opening units, covering 503 ayahs in 29 surahs. Eighty-three cases come from an earlier export and 44 from newer submissions. Event answers are absent. Old ayah splits inherit human review; new production-generated splits were checked and repaired by the preparation assistant, not newly approved by a human. Unresolved cases are excluded. Source spelling is preserved, and references are verbatim canonical clean Quran text.

Starting with 221 candidates outside gold100, we exclude complete gold recitation copies and copies of gold chunks containing annotated events, including such chunks inside longer recordings. We then remove duplicate or redundant ordered training passages. Shared clean ayahs and different error variants remain eligible; this is case separation, not an unseen-ayah split. Gold labels are not exported. Attribution stays private; inputs use surrogate IDs and checksums. Selection targets diversity, not prevalence.

\textbf{Private evaluation.} The fixed human-approved gold100 contains 314 ayah units, including 233 clean units, and 34 opening units: 115 within-ayah and 47 opening events. It covers 274 ayahs in 38 surahs; one of ten learners contributes 46 cases. The selected inputs were previously published in an older task, and twelve also occurred in a machine-labelled pilot. New annotations and membership are private, but this is not an unseen-input or unseen-speaker benchmark.

\textbf{Review material.} The actual training sample contains 23 recordings and 394 units, exceeding ten percent of both combined train-plus-gold recording and unit counts. It contains no selected gold recordings or event answers; shared clean units are allowed. Twenty separate constructed teaching cases, approved by the dataset owner, illustrate the rubric. The train-only interpretation of the workshop's sample requirement needs confirmation, and reviewer access must be finalized. The release is CC BY 4.0 with a no-re-identification condition. No private gold sample is promised.

\section{Evaluation and baselines}
The primary score is label-aware micro event F1: a correct prediction must match both label and location. One-to-one pairing requires the minimum of transcript-span and reference-span intersection-over-union to be at least 0.50. Empty anchors match exactly, or within one word at similarity 0.5, never nonempty spans. Invalid events receive no match credit. Pairing maximizes total span similarity exactly for up to seven events on each side, with a documented greedy fallback above that size. Exact-span F1, per-label F1 and label-agnostic localization are supporting measures.

A secondary application cost is $(rM+F)/U\times100$, where $M$ counts missed mistake events, $F$ false mistake flags and $U$ scored units. We report $r=1$ and $r=2$ rather than fix one rate: none is empirically established, and the ranking is unchanged from 1:1 to 5:1. Predicting nothing wins only below roughly 0.77:1.

\begin{table}[h]\centering\small
\begin{tabular}{lrrr}\toprule
Method & Label F1 & Exact-span F1 & Localization F1\\\midrule
Empty predictions & 0.000 & 0.000 & 0.000\\
Plain normalized diff & 0.531 & 0.505 & 0.839\\
Adapted production components & 0.518 & 0.505 & 0.799\\\bottomrule
\end{tabular}
\caption{Evaluated on all 348 private units with evaluator v2.1. The production baseline adapts cleaner/alignment components; it is not the entire application.}
\end{table}

The diff has zero F1 on eight labels; the adapted components on six. Their high localization relative to label F1 shows why detecting a difference is insufficient. These are task baselines, not results of the new agent comparison.

\section{Competition procedure and feasibility}
Entrants may annotate training data, train models, write rules or use agents, and choose their own development checks. Neither an annotation phase nor an autoresearch loop is required. Organizers rank only the frozen final solution, executed on private inputs with no gold feedback during development. Training annotations may be shared as optional machine-generated artifacts, without implying human validation.

Data, sample, reference, evaluator and a baseline support local preparation; the score script is not a sandbox. Execution isolation, resource rules, calendar and anonymous reviewer access remain to be finalized. A companion paper studies the dataset and evaluator and plans an agent comparison.
\end{document}
